\section{Assessment (revision 2)}\label{sec:assessment}

This section records an assessment of the scheme made on 2026-10-17 against
the maintained implementation~\cite{EigenX}. It is still the authors' own
analysis, not an external review. Its method was as follows. We recomputed
every derived number in Tables~\ref{tab:params}, \ref{tab:attacks} and
\ref{tab:tiers} from the parameters and the cost model, and reread the
proofs of \S\ref{sec:correctness}. We read the key-generation filter that
defines the degree search space. We ran the maintained test suite. We
implemented and measured one new attack component, the test of
Lemma~\ref{lem:pktest}. Records:
\art{crypto:research/eigen_seta_assessment_20261017/}.

\paragraph{Verdict.} We found no attack below the costs of
Table~\ref{tab:attacks}. The binding attack is still mask guessing, at about
$2^\lambda$ classical chains and $2^{\lambda/2}$ Grover queries on the L sets
and $2^\lambda$ queries on L5, Q3 and Q5. Security rests on two assumptions
that have no reduction to a standard problem. Assumption~\ref{ass:key} is that
the hidden-degree eigenbasis with its hidden orientation hides $f$.
Assumption~\ref{ass:msg} is that the masked eigen-torsion problem is hard, and
IND-CPA needs its decisional form. We rank the risks as follows.
\begin{enumerate}
\item \emph{An effective orientation of $\EA$.} The public key is a segment
  of length $2E$ of a horizontal $2$-walk in the hidden
  $\Z[f\omega]$-orientation. Any method that extends this to an evaluable
  class-group action gives subexponential quantum key recovery through
  Kuperberg's algorithm. It does so at every parameter set, because the Q sets
  enlarge $d$ and not $f$. The level-structure
  frameworks~\cite{DFP24,Arpin24} and sesquilinear pairings are the tools to
  check first. Nobody has done that.
\item \emph{A mask-free higher-dimensional embedding.} The subgroups
  $\varphi\langle P_A\rangle$ and $\varphi\langle Q_A\rangle$ and the product
  $ab$ are public and do not depend on the mask ratio. A Kani embedding of
  dimension~4 or~8 built from them alone would remove the binding attack's
  $\sqrt d$ search. This is open.
\item \emph{The decisional assumption.} Every basis of $\EA[2^E]$ is an
  eigenbasis of some endomorphism, so the assumption is about the size of the
  smallest diagonal one. No distinguisher is known, and no argument for
  hardness is known beyond the absence of one.
\end{enumerate}

\paragraph{Corrections to revision 1.}
\begin{itemize}
\item \emph{NIST categories} (relabelling, claim correction). Revision~1
  conjectured categories 1, 3 and 5 for L1, L3 and L5 on the strength of
  AES-like structure. That holds only in the gate-count metric (Tier A of
  \S\ref{sec:nist}). Against $2^\lambda$ quantum queries (Tier B), L1 and L3
  fall short at $2^{64.5}$ and $2^{96.5}$. Only L5, Q3 and Q5 meet that bar.
\item \emph{The degree-search oracle} (advance in the attack's cost
  accounting, small). A candidate degree can be tested on the public key
  alone with a chain of about $2\log_2f$ steps instead of a decapsulation
  of $E\approx4\log_2f$ steps (Lemma~\ref{lem:pktest},
  \S\ref{sec:pktest-results}). This lowers the degree-search row by about
  1.9 bits in the gate model. It changes no tier.
\item \emph{Implementation and measurements} (engineering). The paper now
  points to the maintained code, which decapsulates with a theta chain and
  adds the Q sets. Revision~1's timings and constant-time evidence describe
  the historical Mumford build only (\S\ref{sec:impl-v2}). A field-arithmetic
  defect for non-friendly moduli was found and fixed. It affected no
  parameter set.
\item \emph{Two inconsistencies in the maintained security notes}
  (documentation). \art{eigenx:docs/SECURITY.md} placed the public
  $\omega$-eigenbasis on $\ED[3^F]$. That is impossible, because $3$ is inert
  in $\Q(\sqrt{-7})$, and Proposition~\ref{prop:cyclic} rests on there being
  no eigenline modulo~3. The note is corrected. Its cost model takes the
  admissible fraction of degrees to be $1/\ln2^E$ and says this is the
  density of sums of two squares. Key generation actually accepts $f$ when
  $2^E-7d^2f^2$, after removing prime factors up to $2^{14}$, leaves 1 or a
  prime $\equiv1\pmod4$, together with the divisibility and sign conditions of Algorithm~\ref{alg:keygen}. The measured rate $\rho$
  (\S\ref{sec:key}) is the right quantity. It exceeds the model's fraction
  by 1.6 to 2.3 bits, so the model understates the search, which is the
  conservative direction.
\end{itemize}

\paragraph{Practicality.} Primes of about $10.5\lambda$ bits (1341 bits at
L1, 3904 and 5192 bits for Q3 and Q5) make public keys and ciphertexts 1.3 to
5.2\,kB. Key generation is a rejection search that takes hours at Q5 even
with eight workers (\art{eigenx:research/q5_roundtrip.md}). The decryption
chain is the cost driver. Shortening it with two half-chains (open
problem~3) would bring $p$ from about $d^5$ to about $d^3$, and it is the
most valuable engineering change.
